\documentclass[11pt,a4paper]{article}

% Pakete
\usepackage[utf8]{inputenc}
\usepackage[T1]{fontenc}
\usepackage[ngerman]{babel}
\usepackage[left=18mm,right=18mm,top=18mm,bottom=18mm]{geometry}
\usepackage{helvet}
\renewcommand{\familydefault}{\sfdefault}
\usepackage{amsmath,amssymb}
\usepackage{graphicx}
\usepackage{tikz}
\usetikzlibrary{arrows.meta,positioning,patterns}
\usepackage{tcolorbox}
\tcbuselibrary{skins,breakable}
\usepackage{enumitem}
\usepackage{tabularx}
\usepackage{colortbl}
\usepackage{qrcode}
\usepackage[hidelinks]{hyperref}

% Fachfarbe Physik
\definecolor{physik}{HTML}{2c5aa0}
\definecolor{physiklight}{HTML}{e8f0fa}

% Box-Stile
\tcbset{
    merkbox/.style={
        colback=physiklight,
        colframe=physik,
        fonttitle=\bfseries,
        title=#1,
        boxrule=1pt,
        arc=2pt,
        left=5pt,
        right=5pt,
        top=5pt,
        bottom=5pt
    },
    lueckenbox/.style={
        colback=white,
        colframe=physik,
        boxrule=1pt,
        arc=2pt,
        left=6pt,
        right=6pt,
        top=6pt,
        bottom=6pt
    }
}

% Kopfzeile
\usepackage{fancyhdr}
\pagestyle{fancy}
\fancyhf{}
\lhead{\small Physik 12 GK}
\chead{\small Quantenphysik}
\rhead{\small Name: \underline{\hspace{4cm}}}
\renewcommand{\headrulewidth}{0.4pt}

% Keine Einrückung
\setlength{\parindent}{0pt}
\setlength{\parskip}{6pt}

\begin{document}

% Titel mit QR-Codes
\begin{minipage}[t]{0.65\textwidth}
\vspace{0pt}
{\Large\textbf{Doppelspalt mit Elektronen}}\\[0.3em]
{\normalsize Mi 21.01.2026}
\end{minipage}
\hfill
\begin{minipage}[t]{0.32\textwidth}
\vspace{0pt}
\raggedleft
\begin{tabular}{@{}c@{\hspace{8pt}}c@{}}
\href{https://devbox42.github.io/sciencesim/physik/kl12-GK/01-photoeffekt/04-ES-doppelspalt/SIM-04-doppelspalt.html}{\qrcode[height=1.1cm]{https://devbox42.github.io/sciencesim/physik/kl12-GK/01-photoeffekt/04-ES-doppelspalt/SIM-04-doppelspalt.html}} &
\href{https://devbox42.github.io/sciencesim/physik/kl12-GK/01-photoeffekt/04-ES-doppelspalt/LP-04-doppelspalt.html}{\qrcode[height=1.1cm]{https://devbox42.github.io/sciencesim/physik/kl12-GK/01-photoeffekt/04-ES-doppelspalt/LP-04-doppelspalt.html}} \\
{\scriptsize Simulation} & {\scriptsize Lernpfad}
\end{tabular}
\end{minipage}

\vspace{0.3em}
\hrule height 1pt
\vspace{0.5em}

% ============================================================
% KASTEN 1: VERSUCHSAUFBAU
% ============================================================

\begin{tcolorbox}[merkbox={Kasten 1: Doppelspalt-Experiment}]
\textbf{Aufbau:} Ergänze die Beschriftung in der Skizze.

\vspace{0.5em}
\begin{center}
\begin{tikzpicture}[scale=0.9]
    % Quelle
    \draw[thick, fill=gray!20] (0,0) circle (0.4);
    \node at (0,-0.9) {\small \underline{\hspace{2cm}}};

    % Pfeil
    \draw[->, thick] (0.6,0) -- (2.2,0);

    % Doppelspalt
    \fill[gray!60] (2.5,-1.5) rectangle (2.8,1.5);
    \fill[white] (2.5,-0.6) rectangle (2.8,-0.2);
    \fill[white] (2.5,0.2) rectangle (2.8,0.6);
    \node at (2.65,-2) {\small \underline{\hspace{2cm}}};

    % Schirm
    \fill[gray!30] (6,-1.5) rectangle (6.3,1.5);
    \node at (6.15,-2) {\small \underline{\hspace{2cm}}};

    % Interferenzmuster (angedeutet)
    \foreach \y in {-1.2,-0.6,0,0.6,1.2} {
        \fill[physik!60] (6.3,\y-0.15) rectangle (6.5,\y+0.15);
    }
\end{tikzpicture}
\end{center}

\vspace{0.3em}
\textbf{Prinzip:} Wellen, die durch zwei Spalte gehen, \underline{\hspace{3cm}} miteinander.

An manchen Stellen verstärken sie sich (helle Streifen), an anderen \underline{\hspace{3cm}} sie sich aus (dunkle Streifen).
\end{tcolorbox}

\vspace{0.5em}

% ============================================================
% AUFGABE 1: VORHERSAGE UND BEOBACHTUNG
% ============================================================

\textbf{Aufgabe 1: Vorhersage und Beobachtung} \hfill \textbf{(5 BE)}

\textbf{a)} Was erwartest du auf dem Schirm, wenn \textbf{Elektronen} durch den Doppelspalt geschickt werden? Kreuze an und begründe kurz. \hfill (2 BE)

\vspace{0.3em}
\begin{tabular}{@{}p{0.48\textwidth}p{0.48\textwidth}@{}}
$\square$ Zwei Streifen (wie bei Kugeln) & $\square$ Interferenzmuster (wie bei Wellen) \\
\end{tabular}

\vspace{0.3em}
Begründung: \hrulefill

\hrulefill

\vspace{0.5em}

\textbf{b)} Beschreibe, was du in der Simulation beobachtest, wenn die Intensität immer weiter verringert wird. \hfill (3 BE)

\vspace{2.5cm}

% ============================================================
% AUFGABE 2: DAS RÄTSEL
% ============================================================

\textbf{Aufgabe 2: Das Rätsel der einzelnen Elektronen} \hfill \textbf{(9 BE)}

\textbf{a)} Ergänze den Lückentext. \hfill (4 BE)

\begin{tcolorbox}[lueckenbox]
Jedes einzelne Elektron landet als \underline{\hspace{2cm}} auf dem Schirm. Das spricht für \underline{\hspace{3cm}}.

Aber: Viele Elektronen zusammen bilden ein \underline{\hspace{3cm}}. Das spricht für \underline{\hspace{3cm}}.

Das Besondere: Jedes Elektron geht \underline{\hspace{2cm}} durch den Doppelspalt – es gibt kein anderes Elektron, mit dem es interferieren könnte!
\end{tcolorbox}

\newpage

\textbf{b)} Erkläre in eigenen Worten, warum sich trotzdem ein Interferenzmuster aufbaut. \hfill (5 BE)

\textit{Hinweis: Was „weiß" jedes einzelne Elektron über die beiden Spalte?}

\vspace{3.5cm}

% ============================================================
% KASTEN 2: WELLE-TEILCHEN-DUALISMUS
% ============================================================

\begin{tcolorbox}[merkbox={Kasten 2: Welle-Teilchen-Dualismus}]
Elektronen (und andere Quantenobjekte) sind weder reine \underline{\hspace{2cm}} noch reine \underline{\hspace{2cm}}.

\vspace{0.3em}
Sie zeigen je nach Experiment \underline{\hspace{4cm}} oder \underline{\hspace{4cm}}.

\vspace{0.3em}
\textbf{Faustregel:}
\begin{itemize}[nosep, leftmargin=*]
    \item \textbf{Ausbreitung/Weg:} Wellenverhalten (Interferenz, Beugung)
    \item \textbf{Nachweis/Messung:} Teilchenverhalten (Punkt auf dem Schirm)
\end{itemize}
\end{tcolorbox}

\vspace{0.5em}

% ============================================================
% KASTEN 3: KOMPLEMENTARITÄT
% ============================================================

\begin{tcolorbox}[merkbox={Kasten 3: Komplementaritätsprinzip}]
\textbf{Experiment:} Man stellt einen Detektor auf, um zu messen, durch welchen Spalt das Elektron geht.

\vspace{0.3em}
\textbf{Ergebnis:} Das Interferenzmuster \underline{\hspace{3cm}}!

\vspace{0.5em}
\textbf{Komplementarität (Niels Bohr):}

Wir können \textbf{entweder} den Weg des Elektrons bestimmen (Teilchenbild)

\textbf{oder} Interferenz beobachten (Wellenbild) –

aber \underline{\hspace{4cm}}.
\end{tcolorbox}

\vspace{0.5em}


\vspace{0.5em}

% ============================================================
% ZUSATZAUFGABE (OPTIONAL)
% ============================================================

{\large\textbf{Zusatzaufgabe} ($\bigstar\bigstar\bigstar$)}

\vspace{0.3em}

\textbf{Aufgabe 3:} Ein Mitschüler fragt: „Durch welchen Spalt geht das Elektron denn nun wirklich?"

\textbf{Begründe}, warum diese Frage im Rahmen der Quantenmechanik nicht sinnvoll beantwortet werden kann. \hfill (6 BE)

\textit{Hinweis: Was passiert, wenn man versucht, die Frage experimentell zu beantworten?}

\vspace{4cm}

\begin{center}
\rule{10cm}{0.5pt}

\textbf{Zusatzaufgabe:} \underline{\hspace{1.5cm}} / 6 BE \hspace{2cm} \textbf{Gesamt:} \underline{\hspace{1.5cm}} / 20 BE
\end{center}

\end{document}
